Fix \(t\) and \(z\) outside the \(P_X\)-null exceptional set, and abbreviate the conditional residual by \(\varepsilon\). Conditional centering is part of \(\mathrm{ass:gaussian\text{-}errors}\). For \(u>0\) and \(\lambda>0\), the pointwise bound
\[
 e^{\lambda u}\mathbf 1\{\varepsilon>u\}\le e^{\lambda\varepsilon}
\]
and the conditional moment-generating-function inequality give
\[
 \Pr_P(\varepsilon>u\mid X=z)
 \le\exp\!\left(-\lambda u+\frac{\sigma^2\lambda^2}{2}\right).
\]
Because \(\sigma>0\), the valid choice \(\lambda=u/\sigma^2\) yields
\[
 \Pr_P(\varepsilon>u\mid X=z)\le e^{-u^2/(2\sigma^2)}.
\]
Applying the same argument to \(-\varepsilon\) gives
\[
 \Pr_P(|\varepsilon|>u\mid X=z)\le2e^{-u^2/(2\sigma^2)}. \tag{15}
\]
Pointwise,
\[
 |\varepsilon|=\int_0^\infty\mathbf 1\{|\varepsilon|>u\}\,du.
\]
Conditional Tonelli applied to this nonnegative integrand, followed by the substitution \(v=u/\sigma\), now implies
\[
 \begin{aligned}
 \mathbb E_P[|\varepsilon|\mid X=z]
 &=\int_0^\infty\Pr_P(|\varepsilon|>u\mid X=z)\,du\\
 &\le2\int_0^\infty e^{-u^2/(2\sigma^2)}\,du
 =\sigma\sqrt{2\pi}.
 \end{aligned}\tag{16}
\]
Since \(g_{t,P}(z)\) is finite and \(\mathbb E_P[\varepsilon\mid X=z]=0\),
\[
 \mathbb E_P[Y\mid X=z]
 =g_{t,P}(z)+\mathbb E_P[\varepsilon\mid X=z]
 =g_{t,P}(z).
\]